% ---------------------------------------------------------------
\begin{abstract}
Deployed classifiers are usually monitored through input drift tests, while the labels
that reveal actual harm arrive late. Existing detectors and drift visualizations rarely
test whether an alarm reflects lost accuracy, or when labels arrive. This paper presents \idm{}, a framework that keeps a
label-free input channel (Holm-corrected Kolmogorov--Smirnov tests, a correlation-change
statistic and feature ranking) separate from an impact channel (isotonic-calibrated
average confidence with a Poisson-binomial interval and a harm-concern score). A broker
releases ground truth only after a set delay. The framework was evaluated on 360 controlled synthetic
streams (4.32 million rows), a 1,500-seed replication and three natural streams (UCI
gas sensor, ELEC2, NOAA weather). Marginal shifts were detected in 30/30 seeds. When
only $P(Y\mid X)$ changed, none of 1,205 harmful batches received a concern flag, and
1,140 of them received no input alarm either. After onset, mean estimated accuracy was
0.723 against a realised 0.370. The correlation statistic detected a dependence-only
change in 30/30 seeds versus 16/30 for the KS tests (+46.7 points). Combined with KS,
however, it raised stationary alarms from 0.83 to 2.37 per 24 batches. Recalibration on
released labels cut gas-sensor estimation error by 68.9\% (MAE 0.545 to 0.169), but did
not help on ELEC2 (0.059 to 0.063). Monitor features lowered the Brier score of a harm
forecaster by 27.9\% over a time-only model on known laws. On a held-out law, time alone
kept the highest AUC. Input drift alarms are therefore not reliable harm signals, and
until labels are released some harm cannot be detected from unlabelled evidence.
\end{abstract}
